% !TEX root = ../main.tex
\chapter{Conclusiones}
\label{ch:conclusiones}

Esta entrega recoge un avance del proyecto, no su cierre. Se han
desarrollado los fundamentos del campo y se ha deducido la ecuación
radial de Klein--Gordon masiva en la geometría de Schwarzschild,
incluyendo la coordenada tortuga y el potencial efectivo. Esta
formulación identifica el problema que deberá evolucionarse después.

La parte numérica realizada corresponde a siete evoluciones en espacio
plano. Para los cosenos seleccionados, el método reproduce la solución
exacta y la relación de dispersión. El refinamiento conjunto de espacio
y tiempo muestra órdenes próximos a dos. Estos resultados respaldan
el método en las condiciones concretas de la prueba; no constituyen
una validación de la evolución en Schwarzschild.

La oscilación observada se explica por la ecuación del modo:
su aceleración es opuesta a su amplitud instantánea. La envolvente
permanece constante porque el modelo plano utilizado no contiene
amortiguamiento y su contorno periódico no deja escapar energía neta.
Que el campo cruce por cero no significa que su energía desaparezca.

Queda por definir el problema de evolución en Schwarzschild con sus
datos iniciales, dominio y condiciones de contorno. Después habrá que
ejecutarlo y comprobar su convergencia antes de interpretar cualquier
cambio de amplitud. Las conclusiones físicas de esa etapa se
redactarán a partir de sus resultados.
